With the legal deadline for universal sanitation set for 2033,
municipal rankings such as the ABES Sanitation Universalization
Ranking steer public attention and investment. In this context,
this study investigates to what extent the choice of multi-criteria
method changes the ranking of municipalities in Northeast Brazil by
sanitation, and whether any ranking is more strongly associated
with hospitalizations for diseases related to inadequate environmental
sanitation (DRSAI). To this end, the 505 municipalities eligible for
the 2026 ABES ranking were ranked on its five criteria by crossing
four weighting methods (equal weights, CRITIC, classical AHP, and
Gaussian AHP) with four aggregation strategies (weighted sum, TOPSIS,
VIKOR, and PROMETHEE II). Next, with all parameters fixed in advance,
the 16 rankings were correlated (Spearman) with the DRSAI rate in the
360 municipalities for which this rate was available, using bootstrap
intervals for the differences. As a result, weights changed the
ordering more than aggregation did, and up to 300 municipalities
changed quartile relative to the official ranking. Moreover, in every
ranking, better-placed municipalities had fewer hospitalizations
($\rho$ from $-0.11$ to $-0.23$), with the median rate in the worst
quartile about twice that of the best. In this respect, the official
ranking proved as consistent with health as the alternatives, whereas
seven combinations reduced this consistency. Finally, an exploratory
consensus among the nine rankings that preserved this consistency
identifies 33 municipalities, all lacking adequate waste disposal, as
priorities for public policy under any of these choices.